\documentclass[a4paper]{arrowhead}

\usepackage[yyyymmdd]{datetime}
\usepackage{etoolbox}
\usepackage[utf8]{inputenc}
\usepackage[T1]{fontenc}
\usepackage{multirow}
\usepackage{longtable}
\usepackage{tikz}
\usepackage{adjustbox}
\AtBeginEnvironment{longtable}{\small}
\BeforeBeginEnvironment{tikzpicture}{\begin{adjustbox}{max width=\textwidth}}
\AfterEndEnvironment{tikzpicture}{\end{adjustbox}}
\usetikzlibrary{positioning,shapes.multipart,arrows.meta}

\renewcommand{\dateseparator}{-}
\setlength{\emergencystretch}{3em}
\lstset{language=HTTP,breakatwhitespace=false,columns=fullflexible}

\newcommand{\fparam}[1]{\textit{\textcolor{ArrowheadBlue}{#1}}}

%% Special references
\newcommand{\mref}[2][]{{\textcolor{ArrowheadPurple}{\hyperref[sec:model:#2]{\ifstrempty{#1}{#2}{#1}}}}}
\newcommand{\oref}[1]{{\textcolor{ArrowheadBlue}{\hyperref[sec:operations:#1]{#1}}}}
\newcommand{\pdef}[1]{{\textcolor{ArrowheadGrey}{#1 \label{sec:model:primitives:#1} \label{sec:model:primitives:#1s}}}}
\newcommand{\pref}[1]{{\textcolor{ArrowheadGrey}{\hyperref[sec:model:primitives:#1]{#1}}}}

% as template_idd.tex; the input and output arguments are passed already formatted (\mref or \pref)
\newrobustcmd\osubsection[5]{
  \addtocounter{subsection}{1}
  \addcontentsline{toc}{subsection}{\protect\numberline{\thesubsection}operation \textcolor{ArrowheadBlue}{#1}}
  \subsection*{
    \thesubsection\quad
    #2 \textcolor{ArrowheadPurple}{#3} \\
    \small
    \hspace*{0.075\textwidth}\begin{minipage}{0.1\textwidth}
      \vspace*{1mm}
      Operation: \\
      \notblank{#4}{Input: \\}{}
      \notblank{#5}{Output: \\}{}
    \end{minipage}%
    \begin{minipage}{0.825\textwidth}
      \vspace*{1mm}
      \textcolor{ArrowheadBlue}{#1} \\
      \notblank{#4}{#4 \\}{}
      \notblank{#5}{#5 \\}{}
    \end{minipage}
  }
  \label{sec:operations:#1}
}
\newrobustcmd\msubsection[2]{
  \addtocounter{subsection}{1}
  \addcontentsline{toc}{subsection}{\protect\numberline{\thesubsection}#1 \textcolor{ArrowheadPurple}{#2}}
  \subsection*{\thesubsection\quad#1 \textcolor{ArrowheadPurple}{#2}}
  \label{sec:model:#2}
}
%%

\usepackage{xurl}
\begin{document}
\sloppy

%% Arrowhead Document Properties
\ArrowheadTitle{systemDiscovery MQTT/TLS/JSON}
\ArrowheadServiceID{systemDiscovery}
\ArrowheadType{Interface Design Description}
\ArrowheadTypeShort{IDD}
\ArrowheadVersion{5.3.0}
\ArrowheadDate{2026-11-09}
\ArrowheadAuthor{Jerker Delsing}
\ArrowheadStatus{REVIEW DRAFT}
\ArrowheadContact{jerker.delsing@ltu.se}
\ArrowheadFooter{\href{www.arrowhead.eu}{www.arrowhead.eu}}
\ArrowheadSetup
%%

%% Front Page
\begin{center}
  \vspace*{1cm}
  \huge{\arrowtitle}

  \vspace*{0.2cm}
  \LARGE{\arrowtype}
  \vspace*{1cm}

  \Large{Service ID: \textit{"\arrowid"}}

  \vspace*{\fill}

  % Front Page Image

  \vspace*{1cm}
  \vspace*{\fill}

  % Front Page Abstract
  \begin{abstract}
    This document describes the MQTT protocol with TLS payload security and JSON payload encoding variant of the \texttt{systemDiscovery} service of the ServiceRegistry system.
  \end{abstract}

  \vspace*{1cm}
\end{center}

\newpage
%%

%% Table of Contents
\tableofcontents
\newpage
%%

\section{Documentation edition and implementation scope}
This is the 5.3.0 documentation edition, dated 2026-11-09. It carries forward the system and service contracts from the retained 5.2.2 documentation and identifies the current primary implementation sources. The former \path|documentation/v5.2.2| folder is retained unchanged.

Document version 5.3.0 is distinct from executable, dependency, registered system/service and wire-contract versions. Values still defined by the source (including 5.2.2 artefact versions, 1.0.0 service versions and \path|SysMonitor-AMQP-IDD-5.2.2| payload identifiers) are preserved. This documentation update does not change runtime behaviour or establish a new implementation release. Historical verification statements describe the tests recorded in the inherited baseline; this revision validates document compilation and references.

\section{Overview}
\label{sec:overview}

This document describes the \texttt{systemDiscovery} service interface of the ServiceRegistry system over MQTT, which provides registration and lookup of the application, support and core systems of the local cloud, with their addresses, version, metadata and hosting device.

It's implemented using protocol, encoding as stated in below table:
\begin{table}[ht!]
  \centering
  \begin{tabular}{|l|l|l|}
    \hline
    \rowcolor{gray!33} Profile type & Type & Version \\ \hline
    Transfer protocol & MQTT & 3.1.1 \cite{mqtt311} \\ \hline
    Data encryption & TLS & 1.3 \cite{rfc8446}, 1.2 \cite{rfc5246} \\ \hline
    Encoding & JSON & RFC 8259 \cite{rfc8259} \\ \hline
    Compression & N/A & - \\ \hline
    Semantics & N/A & - \\ \hline
    Ontology & N/A & - \\ \hline
  \end{tabular}
  \caption{Communication and semantics details used for the \texttt{systemDiscovery} service interface}
  \label{tab:comunication_semantics_profile}
\end{table}

This document provides the Interface Design Description IDD to the \textit{systemDiscovery -- Service Description} document \cite{sr-systemDiscovery-sd}.
The operations, their semantics and the data models are the same as in the HTTP/TLS/JSON IDD \cite{sr-systemDiscovery-http-idd}; this document specifies the MQTT binding.

The MQTT interface is offered only when the system runs with \texttt{mqtt.api.enabled=true}; it is then registered with the interface template \texttt{generic\_\allowbreak{}mqtts} (\texttt{generic\_\allowbreak{}mqtt} without TLS), the broker address and port (\texttt{mqtt.\allowbreak{}broker.\allowbreak{}address}, \texttt{mqtt.\allowbreak{}broker.\allowbreak{}port}) and the base topic \texttt{arrowhead/\allowbreak{}serviceregistry/\allowbreak{}system-\allowbreak{}discovery/\allowbreak{}}. In the Arrowhead-test-cloud-java deployment the MQTT interfaces of the core systems are disabled.

\newpage

\section{Service Operations}
\label{sec:operations}

The provider subscribes to one topic per operation, the base topic followed by the operation name:

\noindent\begin{longtable}{| p{5cm} | p{10.4cm} |} \hline
\rowcolor{gray!33} Operation & Request topic \\ \hline \endhead
\texttt{register} & \texttt{arrowhead/\allowbreak{}serviceregistry/\allowbreak{}system-\allowbreak{}discovery/\allowbreak{}register} \\ \hline
\texttt{lookup} & \texttt{arrowhead/\allowbreak{}serviceregistry/\allowbreak{}system-\allowbreak{}discovery/\allowbreak{}lookup} \\ \hline
\texttt{revoke} & \texttt{arrowhead/\allowbreak{}serviceregistry/\allowbreak{}system-\allowbreak{}discovery/\allowbreak{}revoke} \\ \hline
\end{longtable}

A consumer publishes a request envelope (\mref{MqttRequestTemplate}) to the topic and receives the answer (\mref{MqttResponseTemplate}) on its own \texttt{responseTopic}. The \texttt{authentication} field carries the requester's identity according to the provider's authentication policy: \texttt{IDENTITY-TOKEN//<token>} (outsourced or internal policy), \texttt{SYSTEM//<name>} (declared policy) or the requester's X.509 certificate (certificate policy). The \texttt{status} of the response uses the HTTP status code numbers (200, 201, 204, 400, 401, 403, 404, 408, 423, 500, 503); on error the payload is an \mref{ErrorMessage}.

\begin{lstlisting}[caption={Request envelope.}]
{
  "traceId": "a1b2c3",
  "authentication": "IDENTITY-TOKEN//<token>",
  "responseTopic": "arrowhead/response/TemperatureDisplay/a1b2c3",
  "qosRequirement": 1,
  "params": {},
  "payload": {}
}
\end{lstlisting}

\begin{lstlisting}[caption={Response envelope.}]
{
  "status": 200,
  "traceId": "a1b2c3",
  "receiver": "TemperatureDisplay",
  "payload": {}
}
\end{lstlisting}

\osubsection{register}{PUBLISH}{\texttt{arrowhead/\allowbreak{}serviceregistry/\allowbreak{}system-\allowbreak{}discovery/\allowbreak{}register}}{\mref{SystemRegisterRequest}}{\mref{SystemResponse}}

Registers the requesting system (name from its identity) with addresses, version, metadata and an optional device. Returns 201 for a new entry, 200 if an identical entry exists; an existing system with different data is refused (400).

\begin{lstlisting}[caption={Payload of an \oref{register} request.}]
{
  "metadata": {
    "location": "hall-1"
  },
  "version": "1.0.0",
  "addresses": [
    "string"
  ],
  "deviceName": "Raspberry1"
}
\end{lstlisting}

\osubsection{lookup}{PUBLISH}{\texttt{arrowhead/\allowbreak{}serviceregistry/\allowbreak{}system-\allowbreak{}discovery/\allowbreak{}lookup}}{\mref{SystemLookupRequest}}{\mref{SystemListResponse}}

Returns the systems matching all given filters. With \texttt{verbose=\allowbreak{}true} (and \texttt{discovery.\allowbreak{}verbose} enabled) the device details are included.

Optional entries of \texttt{params}: \texttt{verbose}.

\begin{lstlisting}[caption={Payload of an \oref{lookup} request.}]
{
  "systemNames": [
    "TemperatureSensor"
  ],
  "addresses": [
    "string"
  ],
  "addressType": "IPV4",
  "metadataRequirementList": [
    {
      "unit": "celsius"
    }
  ],
  "versions": [
    "1.0.0"
  ],
  "deviceNames": [
    "Raspberry1"
  ]
}
\end{lstlisting}

\osubsection{revoke}{PUBLISH}{\texttt{arrowhead/\allowbreak{}serviceregistry/\allowbreak{}system-\allowbreak{}discovery/\allowbreak{}revoke}}{}{}

Removes the requesting system and its service instances. Returns 200 if removed, 204 if it was not registered.

\newpage

\section{Data Models}
\label{sec:model}

Here, all data objects that can be part of the service calls associated with this service are listed in alphabetic order.
Each \textit{struct} denotes a JSON \pref{Object} with the named fields; fields that are not needed may be omitted.

\msubsection{struct}{MqttRequestTemplate}

The request envelope.

\noindent\begin{tabularx}{\textwidth}{| p{4.25cm} | p{3.5cm} | X |} \hline
\rowcolor{gray!33} Object Field & Value Type & Description \\ \hline
"traceId" & \pref{String} & Identifier chosen by the requester, returned in the response. \\ \hline
"authentication" & \pref{String} & Identity of the requester (see Section \ref{sec:operations}). \\ \hline
"responseTopic" & \pref{String} & Topic on which the response is published. \\ \hline
"qosRequirement" & \pref{Number} & MQTT QoS of the response (0, 1 or 2). \\ \hline
"params" & \pref{Object}$<$\pref{String}$>$ & Optional parameters (the query parameters of the HTTP binding). \\ \hline
"payload" & \pref{Object} & The operation input. \\ \hline
\end{tabularx}

\msubsection{struct}{MqttResponseTemplate}

The response envelope.

\noindent\begin{tabularx}{\textwidth}{| p{4.25cm} | p{3.5cm} | X |} \hline
\rowcolor{gray!33} Object Field & Value Type & Description \\ \hline
"status" & \pref{Number} & Status code. \\ \hline
"traceId" & \pref{String} & The trace identifier of the request. \\ \hline
"receiver" & \pref{String} & Name of the requester. \\ \hline
"payload" & \pref{Object} & The operation output, or an \mref{ErrorMessage}. \\ \hline
\end{tabularx}

\msubsection{struct}{ErrorMessage}




\noindent\begin{tabularx}{\textwidth}{| >{\raggedright\arraybackslash}p{3.9cm} | >{\raggedright\arraybackslash\small}p{5.3cm} | X |} \hline
\rowcolor{gray!33} Object Field & Value Type & Description \\ \hline
"error\allowbreak{}Message" & \pref{String} &  \\ \hline
"error\allowbreak{}Code" & \pref{Number} &  \\ \hline
"exception\allowbreak{}Type" & \pref{ExceptionType} &  \\ \hline
"origin" & \pref{String} &  \\ \hline
\end{tabularx}



\msubsection{struct}{Address}

A network address with its type.


\noindent\begin{tabularx}{\textwidth}{| >{\raggedright\arraybackslash}p{3.9cm} | >{\raggedright\arraybackslash\small}p{5.3cm} | X |} \hline
\rowcolor{gray!33} Object Field & Value Type & Description \\ \hline
"type" & \pref{String} & Address type: \texttt{HOSTNAME}, \texttt{IPV4}, \texttt{IPV6} or \texttt{MAC}. \\ \hline
"address" & \pref{String} & One network address (host name, IPv4, IPv6 or MAC address). \\ \hline
\end{tabularx}



\msubsection{struct}{DeviceResponse}

A registered device.


\noindent\begin{tabularx}{\textwidth}{| >{\raggedright\arraybackslash}p{3.9cm} | >{\raggedright\arraybackslash\small}p{5.3cm} | X |} \hline
\rowcolor{gray!33} Object Field & Value Type & Description \\ \hline
"name" & \pref{String} & Device name. \\ \hline
"metadata" & \pref{Object} & Free-form metadata (JSON object). \\ \hline
"addresses" & \pref{Array}$<$\allowbreak{}\mref{Address}$>$ & Network addresses with their types. \\ \hline
"created\allowbreak{}At" & \pref{DateTime} & UTC time of creation. \\ \hline
"updated\allowbreak{}At" & \pref{DateTime} & UTC time of the last change. \\ \hline
\end{tabularx}



\msubsection{struct}{MetadataRequirement}

A JSON object of metadata requirements: each key is a metadata path (dot notation), each value the required value or an operator object such as \texttt{\{"op": "LESS\_\allowbreak{}THAN",\allowbreak{} "value": 10\}}.



\msubsection{struct}{SystemListResponse}

List of systems.


\noindent\begin{tabularx}{\textwidth}{| >{\raggedright\arraybackslash}p{3.9cm} | >{\raggedright\arraybackslash\small}p{5.3cm} | X |} \hline
\rowcolor{gray!33} Object Field & Value Type & Description \\ \hline
"entries" & \pref{Array}$<$\allowbreak{}\mref{SystemResponse}$>$ & Matching records of the current page. \\ \hline
"count" & \pref{Number} & Total number of matching records (all pages). \\ \hline
\end{tabularx}



\msubsection{struct}{SystemLookupRequest}

Filter for systems; all fields optional, combined with AND.


\noindent\begin{tabularx}{\textwidth}{| >{\raggedright\arraybackslash}p{3.9cm} | >{\raggedright\arraybackslash\small}p{5.3cm} | X |} \hline
\rowcolor{gray!33} Object Field & Value Type & Description \\ \hline
"system\allowbreak{}Names" & \pref{Array}$<$\allowbreak{}\pref{String}$>$ & System names to match. \\ \hline
"addresses" & \pref{Array}$<$\allowbreak{}\pref{String}$>$ & Network addresses of the entity. \\ \hline
"address\allowbreak{}Type" & \pref{String} & Address type filter: \texttt{HOSTNAME}, \texttt{IPV4}, \texttt{IPV6} or \texttt{MAC}. \\ \hline
"metadata\allowbreak{}Requirement\allowbreak{}List" & \pref{Array}$<$\allowbreak{}\mref{MetadataRequirement}$>$ & Metadata requirements; each map is one alternative. \\ \hline
"versions" & \pref{Array}$<$\allowbreak{}\pref{Version}$>$ & System versions to match. \\ \hline
"device\allowbreak{}Names" & \pref{Array}$<$\allowbreak{}\pref{String}$>$ & Device names to match. \\ \hline
\end{tabularx}



\msubsection{struct}{SystemRegisterRequest}

Registers the requesting system; the name is taken from its authenticated identity.


\noindent\begin{tabularx}{\textwidth}{| >{\raggedright\arraybackslash}p{3.9cm} | >{\raggedright\arraybackslash\small}p{5.3cm} | X |} \hline
\rowcolor{gray!33} Object Field & Value Type & Description \\ \hline
"metadata" & \pref{Object} & Free-form metadata (JSON object). \\ \hline
"version" & \pref{Version} & Version of the system. \\ \hline
"addresses" & \pref{Array}$<$\allowbreak{}\pref{String}$>$ & Network addresses of the entity. \\ \hline
"device\allowbreak{}Name" & \pref{String} & Name of the device hosting the system. \\ \hline
\end{tabularx}



\msubsection{struct}{SystemResponse}

A registered system.


\noindent\begin{tabularx}{\textwidth}{| >{\raggedright\arraybackslash}p{3.9cm} | >{\raggedright\arraybackslash\small}p{5.3cm} | X |} \hline
\rowcolor{gray!33} Object Field & Value Type & Description \\ \hline
"name" & \pref{String} & System name. \\ \hline
"metadata" & \pref{Object} & Free-form metadata (JSON object). \\ \hline
"version" & \pref{Version} & Version of the system. \\ \hline
"addresses" & \pref{Array}$<$\allowbreak{}\mref{Address}$>$ & Network addresses with their types. \\ \hline
"device" & \mref{DeviceResponse} & The device hosting the system, if registered. \\ \hline
"created\allowbreak{}At" & \pref{DateTime} & UTC time of creation. \\ \hline
"updated\allowbreak{}At" & \pref{DateTime} & UTC time of the last change. \\ \hline
\end{tabularx}



\subsection{Primitives}
\label{sec:model:primitives}

As all messages are encoded using the JSON format \cite{rfc8259}, the following primitive constructs, part of that standard, become available.

\noindent\begin{tabularx}{\textwidth}{| p{4.6cm} | X |} \hline
\rowcolor{gray!33} Type & Description \\ \hline
\pdef{Array} & A JSON Array \cite{rfc8259}; \pref{Array}$<$A$>$ has elements of type A. \\ \hline
\pdef{DateTime} & JSON String with an ISO 8601 UTC timestamp, e.g. \texttt{2026-10-22T12:00:00Z}. \\ \hline
\pdef{ExceptionType} & One of \texttt{ARROWHEAD}, \texttt{INVALID\_\allowbreak{}PARAMETER}, \texttt{AUTH}, \texttt{FORBIDDEN}, \texttt{DATA\_\allowbreak{}NOT\_\allowbreak{}FOUND}, \texttt{TIMEOUT}, \texttt{LOCKED}, \texttt{INTERNAL\_\allowbreak{}SERVER\_\allowbreak{}ERROR}, \texttt{EXTERNAL\_\allowbreak{}SERVER\_\allowbreak{}ERROR} (JSON String). \\ \hline
\pdef{Number} & JSON Number, an integer. \\ \hline
\pdef{Object} & JSON Object; \pref{Object}$<$A$>$ has string keys and values of type A. \\ \hline
\pdef{String} & JSON String. \\ \hline
\pdef{Version} & JSON String, semantic version \texttt{major.minor.patch}. \\ \hline
\end{tabularx}

\newpage

\section{Go and Rust implementation variants}
\subsection{Go implementation}
Primary system folder: \path|Eclipse-Arrowhead-go/ah5-core-go/serviceregistry|.\par
The system owns its Go module, launcher and build configuration. Use the per-system \path|Makefile| and \path|go.mod|; the containing Go tree is the build source, not a historical Arrowhead copy. Ordinary native roles use independent identities and JSON snapshot state rather than the Java database. Do not assume Java SQL migration, clustering or identical transaction/retry semantics. The registry has its own configuration and optional native AMQP monitor; other roles must be checked against their own registrations.
\paragraph{Listening port.} Default 8443 (server.port in application.properties). Override with server.port, SERVER\_PORT=N or -server.port N; the host is server.address.
The native registry API, outsourced authentication, JSON persistence, optional AMQP telemetry and differences from Java are specified in \path|Eclipse-Arrowhead-go/ah5-core-go/serviceregistry/UPSTREAM-README.md|. It has no MQTT application API according to that profile.
Native design/route evidence: \path|Eclipse-Arrowhead-go/ah5-core-go/serviceregistry/documentation|.\par
For this Java interface, select only a native transport and operation explicitly implemented and advertised by that port. A matching method/path does not establish DTO, error, authorization, retry or persistence equivalence. Java MQTT/AMQP bindings are not automatically provided by a Go or Rust role. Differential contract and security checks remain necessary before substitution.
\subsection{Rust implementation}
Primary system folder: \path|Eclipse-Arrowhead-rust/ah5-core-rust/serviceregistry|.\par
The per-system entry point is built by the \path|arrowhead-core-rust| crate. Build from the primary Rust workspace with the applicable package/binary selection. Documentation links can resolve into the shared core/support catalogue. Independent process identity and state remain required.\par
The ordinary native profile uses Tokio/Axum and rustls with peer-identity checks and independent locked JSON state. Network effects are not distributed database transactions. The precise trust, CRL reload and persistence guarantees are those of the native profile, not the Java implementation.\par
Native design/route evidence: \path|Eclipse-Arrowhead-rust/ah5-core-rust/documentation/ServiceRegistry|.\par
For this Java interface, select only a native transport and operation explicitly implemented and advertised by that port. A matching method/path does not establish DTO, error, authorization, retry or persistence equivalence. Java MQTT/AMQP bindings are not automatically provided by a Go or Rust role. Differential contract and security checks remain necessary before substitution.

\bibliographystyle{IEEEtran}
\bibliography{bibliography}

\newpage

\section{Revision History}
\subsection{Amendments}

\noindent\begin{tabularx}{\textwidth}{| p{1cm} | p{3cm} | p{2cm} | X | p{4cm} |} \hline
\rowcolor{gray!33} No. & Date & Version & Subject of Amendments & Author \\ \hline

1 & 2026-10-22 & 5.2.2 & Document created from the Arrowhead LaTeX templates for the 5.2.2 Java Spring implementation, including the local extensions of this workspace. & Jerker Delsing \\ \hline
2 & 2026-11-09 & 5.3.0 & New documentation edition; primary-source references and Go/Rust implementation scope; 5.2.2 retained. & Documentation maintenance \\ \hline
\end{tabularx}

\subsection{Quality Assurance}

\noindent\begin{tabularx}{\textwidth}{| p{1cm} | p{3cm} | p{2cm} | X |} \hline
\rowcolor{gray!33} No. & Date & Version & Approved by \\ \hline

1 & & \arrowversion  &  \\ \hline

\end{tabularx}

\end{document}
